%% RESUMO NA LINGUA VERNACULA (elemento obrigatorio -- NBR 6028)
%%
%% Regras:
%%   - paragrafo unico, sem recuo, justificado;
%%   - Times New Roman 12, espacamento simples;
%%   - entre 150 e 500 palavras;
%%   - verbo na voz ativa e na terceira pessoa do singular;
%%   - as palavras-chave sao definidas em config/dados.tex.
%%
%% Um bom resumo informa: objetivo geral; principais conceitos e autores da
%% fundamentacao teorica; etapas e acoes da execucao; populacao envolvida e
%% periodo; principais resultados alcancados; conclusoes.

\begin{resumo}
Este espaço é reservado para apresentar o resumo do seu trabalho na língua
vernácula. Ele deve conter entre 150 e 500 palavras, em Times New Roman
tamanho 12, com espaçamento simples e alinhamento justificado. Um bom resumo
deve informar o objetivo geral do trabalho, os principais conceitos e autores
trabalhados na fundamentação teórica, as principais etapas e ações na execução
da proposta, a população envolvida e o período no qual os dados foram obtidos.
Apresente também os principais resultados alcançados e as conclusões a que se
chegou. O resumo deve ser redigido em um único parágrafo, sem recuo de primeira
linha, empregando o verbo na voz ativa e na terceira pessoa do singular. Evite
símbolos, fórmulas, equações e diagramas que não sejam absolutamente
necessários, bem como citações bibliográficas. Em caso de dúvidas, consulte a
norma NBR 6028.
\end{resumo}
